\section{Repeated shadows and the restriction kernel}
\label{sec:shadows}
This section includes the elementary shadow-fiber theorem of the earlier
phase196 manuscript. It makes explicit the rigidity condition in
Theorem~\ref{UR.Restriction}; the old objects being recognizable does not
force the new objects to be determined by them.

Let $P$ be a finite typed incidence system, $J$ a finite set of new types,
and $Q_j$ pairwise disjoint finite sets, disjoint also from $P$. For
$q\in Q_j$ prescribe its \emph{complete old shadow} $s_j(q)\subseteq P$.
Retain the old incidences, put $p*q$ exactly when $p\in s_j(q)$, and make
all objects of distinct new types incident. There are no incidences within
a new type. Write $\Gamma$ for the extension and
$m_j(S)=|s_j^{-1}(S)|$. Let
\[
H=\{h\in K(P):m_j(hS)=m_j(S)\text{ for every }j,S\}.
\]
Absent shadows have multiplicity zero.

\begin{theorem}[The complete shadow kernel]\label{UR.Shadows}
Restriction to $P$ gives the split exact sequence
\[
1\longrightarrow\prod_{j,S}\Sym(s_j^{-1}(S))
\longrightarrow K(\Gamma)\longrightarrow H\longrightarrow1.
\]
Consequently restriction is faithful exactly when every nonempty shadow
fiber has one object. A splitting depends on chosen enumerations of the
fibers; it is not asserted to be canonical.
\end{theorem}
\begin{proof}
A type-preserving correlation preserves $P$, since old and new types are
disjoint. If its restriction is $h$, incidence preservation says that it
maps the fiber over $(j,S)$ bijectively to that over $(j,hS)$. Thus $h\in H$.
When $h=1$, it preserves each fiber. Conversely arbitrary independent
permutations of the fibers preserve old--new incidence and all new--new
incidences. This proves precisely the asserted kernel, not merely a
subgroup contained in it.

For each nonempty fiber choose an enumeration
$e_{j,S}:\{1,\ldots,m_j(S)\}\to s_j^{-1}(S)$. Extend $h\in H$ by
\[
 e_{j,S}(a)\longmapsto e_{j,hS}(a).
\]
Equality of multiplicities makes this well-defined. It preserves
old--new incidence because $p\in S$ if and only if $h(p)\in hS$;
the other incidences are preserved by construction. Its inverse is the
extension of $h^{-1}$. Applying successive extensions keeps the index $a$
unchanged, so these extensions form a homomorphism and split restriction.
Finally a symmetric group is trivial exactly when its set has at most one
element.
\end{proof}

For full correlations one must separately recognize the old object set
and allow any actual permutation of new types. Thus a signature that
includes a new type does not establish pointwise rigidity if a correlation
fixing old objects could still permute those types. The finite examples
in the second application verify this distinction explicitly.

\begin{corollary}[Recognition by unique shadows]\label{UR.UniqueShadows}
Suppose every correlation of an extension $Y$ recognizes an old structure
$X$, and its restriction preserves all reconstructed old incidences and
colors. Suppose every new object is uniquely determined by its complete
old shadow, even across any new types that correlations may permute.
Let $G\leq C(X)$ be exactly the old actions whose forced permutations of
new objects preserve the prescribed incidences and induce a permutation
of the full palette. Then restriction identifies $C(Y)$ with $G$; $K(Y)$
corresponds to the kernel of the resulting palette action.
\end{corollary}
\begin{proof}
Unique shadows give pointwise rigidity, and the definition of $G$ gives
exactly the permissible extensions. Theorem~\ref{UR.Restriction} applies.
\end{proof}

\begin{proposition}[Keeping ordered positions]\label{UR.OrderedShadows}
Let $P$ be a typed incidence system and $T$ a set of distinct ordered
$d$-tuples of old objects, with a prescribed type for every tuple. Form a
new system with a center for each tuple $t$ and a port $(t,i)$ for each
position $1\leq i\leq d$. Give ports distinct position types, disjoint
from the old and center types, and join $(t,i)$ just to its center and the
old object $t_i$. Keep the old incidences. In the subgroup preserving the
old types, center types and each position type, restriction is an
isomorphism onto the old automorphisms preserving the tuple sets by type.
\end{proposition}
\begin{proof}
For a correlation in the stated subgroup, the unique center and old
neighbor of a port show that the center image has entries $h(t_i)$ in
all positions. Since tuples are distinct, its image and all its ports
are forced. Conversely an old automorphism preserving the tuple sets
extends by $t\mapsto(h(t_1),\ldots,h(t_d))$ and $(t,i)\mapsto(h(t),i)$.
Both maps preserve all declared incidences and are inverse under
restriction. Apply Theorem~\ref{UR.Restriction}.
\end{proof}
This is a precise repair of loss of order at the level of a specified
role-preserving subgroup. It does not assert that the entire correlation
group recognizes those roles, that projective dualities extend, or that
the resulting incidence system is a geometry. Those are additional
requirements, not consequences of this proposition.

\begin{theorem}[Rigid arms encode ordered attachments]\label{UR.RigidArms}
Let $X$ be a finite properly colored graph and let $T$ be a family of
pairwise distinct ordered $d$-tuples of its vertices, with $d\ge3$ and
with a color $c(t)$ assigned to each tuple. Add a center for each tuple
and join it to its $i$th entry by a fresh path of length $\ell_i\ge2$,
where the $\ell_i$ are pairwise different. Give centers their prescribed
colors and give internal path vertices separate colors indexed by
$(i,k)$, their arm and distance from the center. All new colors are
disjoint from the old colors. No other new incidences are added.
Suppose every correlation of the resulting graph $Y$ preserves the old
vertex set $V(X)$. Then restriction identifies $C(Y)$ with the subgroup
of $C(X)$ preserving $T$ and its partition into tuple colors. Under this
identification $K(Y)$ is the kernel of the action on the old colors and
tuple colors; all internal path colors are fixed automatically.
\end{theorem}
\begin{proof}
Each component after deleting the old vertices is a subdivided star with
one vertex of degree $d$, its center. Its arm lengths are
$\ell_1-1,\ldots,\ell_d-1$ and are pairwise different. Every isomorphism
between two such components sends the center to the center, preserves
arm indices by their lengths and preserves each distance from the
center. Looking also at the unique old endpoint of each arm therefore
recovers the entire ordered tuple, not merely its unordered boundary.
Distinct tuples have distinct ordered boundaries, so a correlation
fixing all old vertices fixes each component pointwise.

A correlation restricts to $C(X)$ because all old incidences and colors
are retained. Its action sends each tuple to the ordered tuple of its
endpoint images, with one common permutation of the tuple colors.
Conversely an old correlation satisfying these conditions extends by
that rule on centers and by matching indices and distances on arms.
This is an incidence-preserving bijection, and the inverse is obtained
from the inverse old correlation. The internal colors are fixed and
the other colors transform just as asserted. Apply
Theorem~\ref{UR.Restriction}.
\end{proof}
